% 1. Une discussion sur à quel point la station actuel est mauvaise

\subsection{Capteurs actuels}
Les capteutrs actuel se résument à une qualité d'amateurs.
Par exemple, le DHT11 à des températures de fonction entre 0 et 50 degrées celcius, ce qui ne marchera pas l'hivers.
De plus, il lit l'humidité que de 20\% à 90\%, avec une stabilité médiocre entre capteurs. Il ne peut donc pas être utilisé.

Le capteur DSP310 est un bon capteur. Il fonctionne de -40C à 85C, à une précision de $\pm$ 0.5C et $\pm$ 1 hPa.
Il est plutôt fidèle et fiable, et est utilisé dans des situations commerciales. Sa consommation est optimal, à environ 2$\mu$A.
Il est donc correct de continuer à l'utilisé.

Pour ce qui est de l'ensolleillement, le LS06-S fonctionne de -25C à 75C, donc quasiment correct pour nos hivers. Sa consommation
dépend de la luminosité extérieur, autour de quelques mA. Cependant, son gros problème est le manque d'information dans sa datasheet
pour être en mesure de convertir de façon commercial, ses valeurs en unités connues. Ayant donc une fidélité médiocre.
Il fonctionne mieux comme un capteur de lumière que quelquechose qui donne une valeur d'ensolleillement précis.
De plus, il sature facilement.

Le pire reste cependant la station météo sparkfun elle même. C'est un kit éducationel qui n'est pas fait pour une usage commerciale.
En effet, sa composition en plastique est fébrile et ne supporterais pas d'être à l'extérieur. La présence de températures en dessous
de 0C vont causé des bris avec le gel. La précision de l'anémomètre est au meilleur, médiocre. Même si l'unité est à niveau, les cuillères
ne sont pas balancées. La girouette pour la direction du vent à des résistances tellement pas précise qu'il requiert une calibration 
des directions différentes pour chaque unitées. De plus il est possible d'être entre deux résistances, créant un diviseur de tension qui
fausse les valeurs. Pour finir sur cette station, le pluviomètre permet uniquement de mesuré les précipitations qui tombe directment dedans.
Du moment qu'il y a du vent, la précipitation sera fausse. Sans dire qu'il vas gelé l'hivers et accumulé les débris dedans.

Le gros avantage de la station SparkFun est la nature passif du système, donc une consommation moindre.

\subsection{Protocoles actuels}
\subsubsection{Station météo}
Les protocoles actuel sont très robuste dans la station météo prototype.
En effet, les cables sont courts et la station météo est fait pour être mise à l'extérieur, donc un environnement sans grande interférences
comparé à une usine par exemple. Dans le meilleur des mondes, pour simplifié le tout, avoir le même protocol serait meilleur.
I2C permetterais d'utilisé moin de fils. Un mix de SPI et I2C est correct puisque les fonctions I2C peuvent être réutilisé dans le SPI.
Mais sa nature full-duplex n'est pas nécessaire, les capteurs nont pas besoin d'être lut aussi rapidement.
Le meilleur protocol serait I2C, il est bon, fiable et pas chère. UART fonctionne aussi, mais c'est pas nécessaire du tout. Il prend beaucoup
de mémoire dans le microcontrolleur, sans amené grand avantages.

Ceci dit, le choix du capteur est plus important que le protocol, du simple fait que les distances sont courte.

\subsubsection{Protocoles sans fils}
Bluetooth ne peut pas fonctionnée. La distance entre les stations météo et la station de base est potentiellement trop grande. Quelques
centaines de mètres. Sans mentionné la complexité d'avoir plus que 7 station météos pour ce protocol. De plus, Bluetooth requiert des
liscences pour le commerciale.
De plus, en méthode de souscription et notification, la station de base n'a pas grandes accès sur la station météos.

\subsection{Microcontrolleur actuel}
Le ESP32 est un excellent choix, avec son mode basse énergie il peut consommée aussi bas que quelques micro ampères.
Il supporte toutes les protocoles nécessaires, à assez de GPIO pour répondre à tout les besoins, et est relativement peu couteux
lorsqu'il est implémenté sur un PCB fait sur mesure plutôt qu'un kit de développement.